\renewcommand{\footerfilename}{RemedialUnix.tex}


\section{Remedial Unix}  \label{RemedialUnix}IRAF/PyRAF is considered a ``terminal based command line
environment'', a very flexible way to use small programs to achieve
large goals. GUI environments restrict/channels users into beguiling
'canned' logic. Many people are uncomfortable with command line
work. It requires learning, practice, and frequent use to maintain a
fine edge needed for rigorous scientific data reduction.

This \dhl{Bon Mots} document represents many tricks that one may use
to avoid editing temporary files by hand and to provide an audit trail
for projects.

Go to Section \ref{sec:complexexample} and briefly ponder the code with an
eye towards: ``What was this person thinking!''.

\section*{Python Virtual Environments}

Add a bash function (or equivalent) to your .bash\_rc file.

The file should read something like:

\begingroup \fontsize{10pt}{10pt}
\selectfont
%%\begin{Verbatim} [commandchars=\\\{\}]
\begin{verbatim} 
function pyraf { source $HOME/envs/pyraf/bin/activate ;
                 export PYTHONSTARTUP="$HOME/envs/pyraf/.pythonrc.py"
                 export PYTHONPATH="$HOME/.iraf:$VIRTUAL_ENV/lib/python3.12/site-packages";
                 pyraf -s $*;
                }
\end{verbatim}
\endgroup
%% \end{Verbatim}

This function: 1) gets the virtual environment going for the window;
2) allows a few snippet file (regular python) to be included on
startup; 3) sets the PYTHONPATH to match the venv; and 4) launches the
pyraf package.

My .pythonrc.py file:
\begingroup \fontsize{10pt}{10pt}
\selectfont
%%\begin{Verbatim} [commandchars=\\\{\}]
\begin{verbatim} 
from astropy.io import fits           # open fits files in python
import numpy as np                    # numpy
import pandas as pd                   # pandas
from pyraflogin3 import *             # my .iraf/pyraflogin3.py custom funcs
\end{verbatim}
\endgroup
%% \end{Verbatim}


Python requires virtual environments. Virtual environmnents are bad,
but necessary.  PyRAF, as distributed has Python 3.12 baked
in. Extended use of the Pythonic aspects of PyRAF require lots of
special astropy main and affiliate packages and requires versions to
be somewhat matched. Hence the Python Virtual environment.

Here we use \$HOME/envs/pyraf as the 'venv.

\section{The Unix and the Data ``Big Picture''}

This document is centered around the Ubuntu 24.04 LTS and MacOS
Sequoia and later releases using PyRF 2.17 in a docker container.  It
offers a paradigm developed over years of working with IRAF. This
leverages Unix. The astronomy community uses Apple Mac computers that
still retain enough of the Unix nature to allow easy porting of IRAF
and friends to that platform, even the Apple M1,2,3 and 4 processors.

Win1X requires a VirtualBox (the Microsoft HyperV may not work well
enough). Note: there are some AMD processors\index{AMD!processor
  errors} that will not support virtualization. The client operating
system installed in the Virtual box may be Fedora (the community
version of Red Hat Enterprise Linux or RHEL). This allows installing
ESO packages that leverage the RHEL environment.  \textbf{Note}: Some
AMD processors will not run hypervisors.  \index(RHEL!ESO)
\index{RHEL!Win1X} \index(VirtualBox)

\subsection{Unix Commands and Philosophy}

Unix processes command input one line at a time. Long lines can be
created by appending a \\\\ at the end second and subsequent lines
up to but not including the last line.

The command line is breaks lines into an array of strings, each 
array element is a part of the command. The string is split over
whitespace (spaces and tabs). Spaces can be included by quoting
strings with either the single-quote or the double quote. 
The first entry in the array is the command name itself, and
any following entries are the switches or arguments for that
command. 

Takeaway: Spaces are very inportant in Unix commmands.

A Unix command follows the philosophy that says programs should be
small, perform a specific tight duty and do it well, permit some fine
tuning options. A program takes input from ``standard-in''
(\dhl{stdin}), does its work and puts its output to ``standard-out''
(\dhl{stdout}). Any errrors are reported via ``standard-error''
(\dhl{stderr}). Unix uses a trick called a \dhl{pipe}, represented
by the vertical bar \llbox{|} to send the
output (stdout) from one program into the input (stdin) of the next.  Thus
programs can be linked together on one line. \index{Unix!programs}

There are even more complicated ways to manage the flow of commands that
involve conditionals (and and or), redirecting output etc. Find a good
tutorial on Unix and spend time with the details. This will pay huge
dividends later on.

PyRAF/CL uses Stdout (Note: Case is important!). 

OK, some programs are huge, and this paradigm does not necessarily
apply. But all the basic Unix utility programs like \dhl{ls}, \dhl{find},
\dhl{grep} etc work this way.

PyRAF tasks are \dhl{sticky} -- they maintain a \dhl{.param} file somewhere.
The directory {\textasciitilde{}/.iraf/\dhl{uparams}} stores these, together
with python helpers. 

PyRAF uses the \dhl{unlearn <task>} to 'reset' the lparms to their original 
values.\index{PyRAF!unlearn}.

\subsection*{Psychology}

Changing from a GUI paradigm to a command line paradigm is one of the
hardest things for many people reading these notes to do.

The Unix philosophy is set forth by Ken Thompson, as documented by Doug
McIlroy \cite{McIlroy1978,McIlloryPhilosophy}:

\subsection{File System Structure}

A machine sees its entire collection of disks and other data storage;
USB ports, TTY terminals etc as a ``file''. The ``root'' of the
filesystem starts at "/". Everything hangs off that one place.
It is possible to mount foreign (not of this machine) directories
simply as a directory somewhere under root (/) of this machine.
Therefore no distinction of the disk needs be made.

Linux files do not have a true name! They consist of inodes. Each inode
starts out with a 'filename'. This pair is stored in a special file
call a directory as a 'directory entry'. A 'hard link' exists
for the 'first' instance of the file. It is possible to create
another (additional) hard link to the file -- a responsibility
count is upped by one. If I were to delete the first file, it
would be removed from its original directory but it will live
on, perhaps with a different name, in the other hard-link position.
Deleting a file decrements the responsibility count and if the
count goes to zero -- it is gone for ever.

DELETES CAN NOT BE UNDONE (well easily).

A soft-link can be created. Windows calls the idea a 'shortcut' -- but
in reality a shortcut is a separate file if a special type, that has to
be opened. Using shortcuts in Windows incurs a heavy overhead, where
soft-links in Unix have much lighter overheads.

\begin{quote}
    Make each program do one thing well. To do a new job, build afresh
    rather than complicate old programs by adding new "features". Expect
    the output of every program to become the input to another, as yet
    unknown, program. -- McIlroy
\end{quote}

GIUs are monolithic in design.\footnote{A GUI is a great way to encapsulate
a process where variation in the process is shuned and discouraged. Scientific
data is noisy and physics is opinionated -- this requires a careful attention
to detail not managable with GUI environments.}

One negotiates with Linux's root process to connect to the Linux
system, where the login process drops the user into a
\dhl{shell}. There are several shells to choose from -- here we will
stick with \dhl{bash} \index{bash}. \textbf{\emph{Note:}} On Debian
systems \dhl{/usr/sh} may default to an abbreviated shell called
\dhl{dash}. \index{Unix Shell!intro}

\begin{quote}
``On a UNIX system, everything is a file; if something is not a file, it is a process.'' --
Machtelt Garrels, ``Introduction To Linux: A Hands On Guide''
\end{quote}

Peter Freeman \cite{freeman1975software} offered a model of a computer
(and a good operating system) as one of a \dhl{Processor},
\dhl{Memory} and \dhl{Transducer} (PMT) \index{computer!PMT}. The
processors today have less than one core to multiple cores, symmetric
arrays of cores and massively parallel cores.  Even the core of yore
is being supplanted by compute engines, the current popular one is the
GPU with $10^{\$}$ tiny cores.

\subsection*{Unix Filesystem}

First off: MS Windows does not understand the Unix concepts of
\dhl{hard links} or \dhl{soft links}. They have a \dhl{shortcut}
that is a separate file, that requires two costly open operations.
Windows was introduced when the DOS BIOS understood separate devices
annotated as A:, B:, C: etc. Unix uses \dhl{mounts}, where a device
appears in a file system relative to one and only one starting point,
called \dhl{root}. \index{MS Windows!files}

Unix files are specialized and fall into classifications like: directories,
special files, links, sockets, named pipes.

There is only one file system and it starts at ``/'' referred to as
\dhl{root}.  We tend to think of files as residing on a disk
drive. Multiple drives, and their multiple partitions are simply
``mounted'' at a location under the ``/'' root directory. Examples are
usually called something like /usr2. Disks on other machines are
usually found, conceptually as a sub-directory the /mnt directory.
remote network drives may be mounted relative to root as well. 
Thus:
\vspace{-.15cm}
\begin{enumerate}\addtolength{\itemsep}{-0.5\baselineskip}
%\setcounter{enumi}{N}
   \item   All filenames appear to be the same!
   \item   Files may be located anywhere on an intranet, or specifically the internet.
\end{enumerate}

\subsection*{The Shell ,``commands'', and Scripts}

The computer sees typed text as groups of characters separated by 
\dhl{whitespace}. whitespace is one or more space or tab characters.

This is why we do not like spaces in filenames! PyRAF does not support
spaces in filenames.

Indeed, all of programming takes a stream of complex symbols that, when
taken together, direct the computer to take some safe action. Linux
uses whitespace to set each symbol apart. 

Commands are designed to be typed at the keyboard and are usually
highly abbreviated. 

One of the handiest commands is \dhl{man}.

\example Use the ``\dhl{man}'' command \\
\prose Hey, what does the man (or ls or ln) command do?\\
\exercise Enter ``\dhl{man man}'' at the shell prompt.\\
\takeaway Read about how man works!\\
\exercise[-b] Enter ``\dhl{man ls}'' at the shell prompt.\\
\takeaway Note: it says ``list directory contents'' so the abbreviation
of ``ls'' makes sense. It also offers up a myriad of options.\\
\exercise[-c] Try ``\dhl{ls -l}'' to see details of each file, one per line.\\
\takeaway A way to determine if files are normal files, directories or
links.

One of the oldest Unix commands is ``ls'' used to 'LiSt' directories and
their contents in interesting ways.

\section{IRAF and  and the Data ``Big Picture''}

IRAF is not a program, it is a collection of programs. IRAF has
``programs'' called \dhl{cl} and \dhl{ecl} that behave like
shells. The various tasks may be baked into cl/ecl but usually they are
either other cl scripts of compiled executables. PyRAF is the latest
incarnation of the cl/ecl evolution. It did not make sense to make an
\dhl{eecl} program!

IRAF refers to programs as 'tasks'. They may be in the 'ecl' syntax,
or the extenced cl sytax may be transliterated into FORTRAN, compiled
into ELF files and appear as \dhl{task.x} (.x extension).


\section{Files and Directory Structure}

IRAF/PyRAF 2.18 is installed at a system level, and customized for each user.
The \dhl{mkiraf} command will create a {\textasciitilde}/.iraf directory.

IRAF uses two key files: \dhl{login.cl} and \dhl{loginuser.cl} to select from a myrid
of options and packages one will use for the tasks at hand. The \dhl{login.cl}
file now resides at \dhl{/etc/iraf/login.cl}. The loginuser.cl is located
at {\textasciitilde}/.iraf.loginuser.cl -- this is the file you change.
\textbf{\emph{Note}}: a single line with the word ``keep'' is required at the bottom
of this file.

The loginuser.cl file may be used to add your own special on-off commands,
and to leverage \dhl{foreign} tasks (usually other handy programs on the
path). \index{IRAF!foreign tasks}

Each IRAF task has a set of parameters, managed in the {\textasciitilde}/.iraf/uparam
directory. 

\textbf{BEWARE}: IRAF tasks are ``\dhl{sticky}''. When you change something it remembers the new
thing. The IRAF command \dhl{unlearn <task>} will cause the task to revert to
its defaults. \index{IRAF!sticky}

Other files of note/interest:

\vspace{-.15cm}
\begin{enumerate}\addtolength{\itemsep}{-0.5\baselineskip}
%\setcounter{enumi}{N}
   \item  {\textasciitilde}/.bashrc, include your own {\textasciitilde}/.iraf-aliases
   \item  login.cl  - traditional configuration. \dhl{NEVER} edit this one.
   \item  loginuser.cl - included by login.cl with configuration changes particular to
          users special taste.
   \item  pyraflogin.py - while not official may contain handy python functions
   \item  home\$/uparm - where the system keeps parameter files.
\end{enumerate}
Historically programs were smaller as were the computers. System
administrators added files to /usr directory to be system wide
for that machine's users. There were few libraries, usually statically
linked. 

The Filesystem Hierarchy Standard (FHS) was put forward in 2015, and
specifies that packages like IRAF should be installed under /opt/iraf/...
deep. This will require a modest amount of rework to the existing code.

\section{Getting Started}

Configure the system. Copy the login.cl to ~/.iraf. Copy loginuser.cl script
to ~/iraf. \textbf{\emph{Note}} One has a hidden-file dot.

Add the .iraf.aliases file to \dhl{\$HOME} directory and if you like, add it to
the \dhl{\$HOME/.bashrc} file.

Start PyRAF in one window and start Editor in another window. I
recommend emacs or Sublime Text as editors -- they are GUI
oriented. Remember all things have a learning curve.

Why we use the vi editor. Vi is the ``\dhl{Vi}sual Editor'' that sits
on top of a powerful albeit basic Unix editor ``\dhl{ed}. The \dhl{ed}
editor is the basis for the \dhl{sed} the Unix stream editor. Regular
expressions (think wildcards) are very sophisticated and turn up all
over the Linux/Unix system. They are worth the effort to learn. The
``re'' of the grep program (\dhl{g}lobal \dhl{re}gular expression
\dhl{p}rint) program obviously makes heavy use of regular expressions.

The very basic Linux commands are:

\begin{table}[h!]
\centering
\begin{tabular}{| l | l |}
\hline
Command  & Basic Action   \\
\hline
cd    & change directory    \\ 
ls    & list files     \\ 
pwd   & print the current working directory    \\ 
cat   & \dhl{concatenate} file(s)    \\ 
less  & view content of a file    \\ 
man   & the Manual -- print details of a command \\ 
which & show where on PATH an executable is \\ 
rm    & remove files and directories. \\
%% ones-based: \cline{a-b}
\hline
%%\DeleteShortVerb{|}
\end{tabular}
%%\end{minipage}    %% for footnotes  r@{.}l 
\caption{Very Basic Unix Commands}
\label{table:VeryBasicUnixCommands}
%%} % end small etc
\end{table}


\subsection{Very Basic PyRAF commands}


PyRAF switches: \index{PyRAF!switches}
{\color{darkred}
\begingroup \fontsize{10pt}{10pt}
\selectfont
%%\begin{Verbatim} [commandchars=\\\{\}]
\begin{verbatim} 
-h         List the available options. There are long versions of some options
           (e.g., --help instead of -h) which are also described.

-c command Command passed in as string (any valid PyRAF command)

-e         Turn on ECL mode

-m         Run the PyRAF command line interpreter to provide extra capabilities (default)

-i         Do not start the special PyRAF interpreter; just run a standard
           Python interactive session

-n         No splash screen during startup

-s         Silent initialization (does not print startup messages)

-v         Set verbosity (repeated ’v’ increases the level; mainly useful for 
           system debugging)

-x         No graphics will be attempted/loaded during session

-y         Run the IPython shell instead of the normal PyRAF command shell
           (only if IPython is installed)
\end{verbatim}
\endgroup
%% \end{Verbatim}
}

\begin{table}[h!]
\centering
\begin{tabular}{| l | l |}
\hline
Command  & Use   \\
\hline
cl            & run the cl/wcl editor                        \\
epar          & edit parameters/launch task                  \\
nhedit        & fairly compelete header editor               \\ 
imhead        & list the header values for a file            \\ 
hselect       & select values from headers                   \\ 
imexamine     & all manor of ways to view images             \\ 
apall         & reduce a spectrum image                      \\ 
identify      & identify spectral lines assign wavelenghts   \\
dispcor       & apply identify results to science spectrum   \\
splot         & view/analyze reduced spectrum                \\ 
imstat        & print mean,min,max,std etc for set of images \\ 
%% ones-based: \cline{a-b}
\hline
%%\DeleteShortVerb{|}
\end{tabular}
%%\end{minipage}    %% for footnotes  r@{.}l 
\caption{Very Basic IRAF/PyRAF commands}
\label{table:VeryBasicIRAF/PyRAFcommands}
%%} % end small etc
\end{table}


\section{Working on subsets of files.}

IRAF uses ``list files'' denoted \dhl{@list.txt}. The \llbox{@} key
as the first character of the filename tells IRAF to open that
file and take from it filenames -- one filename per line. This is
real handy.

Some documents have you make lists, then use an editor to change the
list as you go. They recommend filenames like mylist.txt. This is
cumbersome to type. Try using a simple name like \dhl{l.l}. Note the
\llbox{l} key is very near the \llbox{.} key! Handy to type.  If you
see the file \dhl{l.l} you know you can simply remove it.  The \llbox{!}
as the first character on a command line sends the rest of the line
to the local bash script. This escape mechanism allows you to use
all the power of Unix inside cl scripts.

The command \dhl{!ls -1 *Dar* > l.l} takes all filenames with Dar
in it and lists them out in lexicographical (sorted) order. Here
are some example of ways to build list-files.

{\color{verbcolor}
\begingroup \fontsize{10pt}{10pt}
\selectfont
%%\begin{Verbatim} [commandchars=\\\{\}]
\begin{verbatim}
!ls -1 *fits > l.l
cl < ~/iraf/crall.cl
# all the files are now c_<filename>

hselect c_*fits $I "IMAGETYP ?= 'Bias'"  > l.l
hedit @l.l IMAGETYP zero add+ show- ver- update+
!if test -e zmaster.fits; then rm zmaster.fits; fi  # bash to rm file if exists
imcombine @l.l zmaster.fits combine=median

hselect c_*fits $I IMAGETYP ?= 'Light'"  > l.l
hedit @l.l IMAGETYP object add+ show- ver- update+
imarith @l.l - zmaster.fits z_@l.l
imarith z_@l.l - dmaster.fits d_z_@l.l
imarith d_z_@l.l / flatmaster_SII.fits f_d_z_@l.l

# files are now {\color{verbcolor}{\verb#f_d_z_c_filename.fits#}}
\end{verbatim}
\endgroup
%% \end{Verbatim}
}
This makes use of the fact that we can pile on prefixes to
processed files:

\begin{itemize}
\addtolength{\itemsep}{-0.5\baselineskip}
   \item   t - trimmed
   \item   c - cosmic ray reduced
   \item   z - zero subtracted
   \item   d - dark corrected
   \item   f - flat normalized
\end{itemize}

At the end, a flat-normalized file called {\color{verbcolor}{\verb#science.fits#}} will
have the name {\color{verbcolor}{\verb#f_d_z_c_t_science.fits#}}.

Lots of intermediate files lying around: Clean them up!

{\color{verbcolor}
\begingroup \fontsize{10pt}{10pt}
\selectfont
%%\begin{Verbatim} [commandchars=\\\{\}]
\begin{verbatim}
print "Remove residual files."
! (export LC_ALL=C; rm [a-z]_*fits 2> /dev/null;)
\end{verbatim}
\endgroup
%% \end{Verbatim}
}

\textbf{\emph{Prose:}} \textbf{\emph{Hey}} PyRAF send the line to bash; but! using a sub-shell
set {\color{verbcolor}{\verb#export#}} the {\color{verbcolor}{\verb#LC_ALL=C;#}}
shell variable to let bash actually use case sensitive wildcards. Then
{\color{verbcolor}{\verb#rm [a-z]_*fits 2> /dev/null;#}} remove the
file -- and in case there is some whining about the process, send the
whines to the bit-bucket ({\color{verbcolor}{\verb#/dev/null#}}).


\section{Observing}
PyRAF \index{PyRAF!Observing} is a useful tool for spectroscopy. Using
the imexamine command and the 'j' key: set the rplot to wider than a bright
line width. 

% HEREHEREHERE

\clearpage
\section{Overall Philosophical Approach}

This document covers a philosophy --- a set of general tools and
scripts to take and reduce data.

In this case some location like {\color{verbcolor}{\verb#/opt/myobs/iraf/#}}
or {\color{verbcolor}{\verb#~/iraf/#}} contains a number of small generic
{\color{verbcolor}{\verb#.cl#}} scripts to handle specific needs. These
are specific to certain conditions -- like a package that writes {\color{verbcolor}{\verb#.fts#}}
instead of {\color{verbcolor}{\verb#.fits#}} files. Routines to remove pesky
spaces that GUI users like to use; get rid of multiple ``dots''; plus signs
etc. 

\vspace{-.15cm}
\begin{enumerate}\addtolength{\itemsep}{-0.5\baselineskip}
   \item   Create a planning directory. Rename to Observations/directory
when the plan is executed.
   \item   Create a ``usw'' directory to hold the plan, collect pre-analysis information:
\vspace{-.15cm}
\begin{enumerate}\addtolength{\itemsep}{-0.5\baselineskip}
   \item   DSS image -- ds9 Analysis \menu Image Servers
   \item   Region files, catalog files for ds9 -- ds9 Region \menu Shape ...
   \item   Photometry reference stars --- ds9 Analysis \menu Catalogs \menu Database \menu UCAC4;
or with the Aladin application and TOPCAT make one the hard way. (PanSTARRS DR1).
   \item   Bibliographic information -- Browser and SIMBAD for the target(s).
   \item   Gather any existing bad pixel masks -- previous runs; observatory supplied files.
\end{enumerate}

   \item  Gather raw data --- stay up all night
   \item  Archive raw data as a RawData and make read only --- back to your {\color{verbcolor}{\verb#Observations#}} directory:\\
{\color{verbcolor}{\verb#(cd Observations/YYYYmmDD; mkdir -p RawData; cd RawData; rsync <from somewhere> .; chmod -w -R *#}}).
\end{enumerate}


\vspace{-.15cm}
\begin{enumerate}[resume]\addtolength{\itemsep}{-0.5\baselineskip}
   \item   Copy raw data to a PreAnalysis directory -- do not damage any
raw data whatsoever. Make the PreAnalysis files writable:\\
({\color{verbcolor}{\verb#cd Observations/YYYYmmDD; mkdir -p PreAnalysis; cd PreAnalysis; cp -pr ../RawData/* .; chmod +w -R .#}})\\
Clean data and make initial analysis files:
\vspace{-.15cm}
\begin{enumerate}\addtolength{\itemsep}{-0.5\baselineskip}
   \item   Fix the filenames --- {\color{verbcolor}{\verb#cl < ~sbo/iraf/fix_sbig.cl#}}
   \item   Fix the headers--- {\color{verbcolor}{\verb#cl < ~sbo/iraf/fix_headers.cl#}}
   \item   Trim (remember) overscan --- see above
   \item   Remove bad WCS especially from zero/dark/flats {\color{verbcolor}{\verb#cl < ~sbo/iraf/fix_sbig_wcs.cl#}}
   \item   Remove cosmic rays ---  {\color{verbcolor}{\verb#!ls -1 *fits > l.l#}}
followed by {\color{verbcolor}{\verb#cl < ~sbo/iraf/crall.cl#}}
\item Other basic reduction steps to be covered elsewhere:
\begin{enumerate}\addtolength{\itemsep}{-0.5\baselineskip}
   \item   Make any bad pixel masks 
   \item   Make master zeros and darks 
   \item   Make master flats 
   \item   apply zeros,darks,flats to science images
   \item   add new WCS or de-distort and add new WCS
   \item   extract photometry
   \item   combine images for deeper photometry
   \item   freeze this work
\end{enumerate}
\end{enumerate}


   \item  Copy the PreAnalysis data into an Analysis Directory.
(Any damage here is quickly recovered from PreAnalysis.)\\
{\color{verbcolor}{\verb#!(cd ..; mkdir -p Analysis; cd Analysis; cp -pr ../PreAnalysis/f_d_z_c_t* .)#}}
\end{enumerate}


A break down of the above steps reveals the need to make/use
certain scripts (some shown in red in context above):

\vspace{-.15cm}
\begin{enumerate}[resume]\addtolength{\itemsep}{-0.5\baselineskip}
   \item   Fix the filenames: remove dots, spaces, dashes other special characters.
   \item   Fix the headers: Observatory/Vendor non-IRAF keywords to IRAF
keywords; Very-Non-Standard keywords like PCOUNT and GCOUNT
   \item   Pinpoint is good enough for tracking (still takes a while to apply
and will fail) but is not accurate.
   \item   Single Pixel errors from cheap cameras, inevitable cosmic rays need
to be mitigated in raw zero, dark and flat frames.
   \item   Masks provide locations in pixel coordinates of defects that can ruin
the science.
   \item   Adding a new WCS will help to quickly align images. This may be enough
if there is no distortion, de-distort with flux conservation where possible.
\end{enumerate}


You can stop here and defer photometry to the analysis step. However,
you can extract photometry using sextractor and a few tricks.

Combining images is an art.


\textbf{\emph{Main goal:}} to develop a \textbf{\emph{compact}} set of
skills with Unix/Python/PyRAF/IRAF to create a
\textbf{\emph{collection}} of IRAF scripts to process images from two
cameras, three telescopes using vendor's software. In this sense,
``compact'' means: the basic IRAF commands; python, bash and Unix
tricks and commands; a basic use of SAOImage/ds9. A familiarity with
command line operations and vocabulary to support planning, executing,
reducing data and publishing photometic data. Summary: the course is
designed to introduce astrophysical hands-on research -- not the
tedious nature of the related software tools.

Examples, long and short, will provide an overview of all the weird/odd ways
Bash/IRAF/PyRAF and Python are used is in Appendix \ref{app:PutItTogether}.

\textbf{\emph{The Problem:}} PyRAF is a ``pythonic''
wrapper/controller of collection of IRAF tasks/commands that in turn
are based on the Unix operating system. The long path to getting a
handle on PyRAF is to build on knowledge gained over time: starting
with Python, then Unix in general and the bash shell in particular,
then IRAF commands and tasks. This is hard to do in two weeks.

\textbf{\emph{Exercise:}} we'll jump into the middle of the mess and
build as we go. To temper this exercise, examples are from a third
observatory and vendor's software -- this is to
\textbf{\emph{prevent}} a simple cut/paste without taking the time to
understand each fragment of the provided example code.



\section{PyRAF, IRAF and Unix (bash)}

IRAF was developed in a slow rather ad hoc fashion, its birthday
established by Donald Wells \cite{FITSBirthday} as 28 March 1979,
real development started in 1981 and was it put to use in 1984
\cite{1993ASPC...52..173T}.  At the time researchers at various
institutions using various computing platforms developed ``packages''
that were brought together at NOAO under IRAF (Image Reduction and
Analysis Facility). It was developed under Digital Equipment
Corporation's PDP and Vax computing environments as well as Unix. The
result today reflects coding and use patterns in vogue since that
time. \index{IRAF!history}

The original IRAF used code from a book that ran afoul of copyright
infringements. While still suited for academic work, the work
was dropped by AURA and taken up by a Github based community. There
were several issues with the older 2.14 code. FORTRAN's COMMON and
EQUIVALENCE blocks did not translate well from 32 bit architectures
to 64 bit architectures. The copyright and FORTRAN issues were resolved,
and released as version 2.17.

\clearpage
From \cite{NOIRLabReleaseNote} page:

\begin{quote}
The Image Reduction and Analysis Facility (IRAF) is a general-purpose
software system for the reduction and analysis of scientific data. Its
development started in 1984 at the National Optical Astronomy
Observatories (NOAO) in Tucson, Arizona. As of January 2024, the
Community Science and Data Center (CSDC) and the US National Gemini
Office (US NGO) at NSF NOIRLab launched the new NOIRLab IRAF v2.18.
\end{quote}

For many years we've been awaiting those new packages to be written.
The amount of legacy code, the need to audit results, should see
IRAF supported well into the future.

PyRAF is a ``pythonic'' wrapper for a vast array of IRAF commands and
``tasks''.  IRAF's early control program was ``cl'' melding ideas from
many system's various approaches.  \index{IRAF!cl!history} Users
wanted ``\emph{more}'' so the extended command language or ``ecl'' was
created. Users wanted more `\emph{`more}''.  So, around 1998, rather
than create an ``extended extended command language'', the
\index{PyRAF!history} IRAF developers adapted the then-current
interactive python base to create PyRAF (\cite{2006hstc.conf..437G}
and references therein). PyRAF is a decent replacement for IRAF/ecl
that is mostly backward compatible (to a very large degree -- and
there are differences) with ecl. In cl/ecl it was always possible to
send a string to the underlying operating system by starting the line
with an exclamation point ({\color{verbcolor}{\verb#!#}}).

UNIX\texttrademark\; (Unix) was licensed to outside parties in the 1970's.
\index{Unix!history}. The philosophy was founded on the premise that things
should be of a minimalist and modular set of clear and concise processing
steps. This is in stark contrast to the modern (haha) use of a Graphical
User Interface (GUI). The GUI places severe restrictions on one's ability
to combine steps in interesting ways to push the bounds of knowledge. It
also acts as a diaper to prevent people from messing up their business
environment. A good tool in one context (business) and very bad in science.

Under Unix, there are about 160 user commands. There is similar number
of terminal actions within a GUI so learning commands is not that daunting.

Summaries of basic Unix commands abound on the net. Here we will
ignore the basics (ls,cp,mv,rm,pwd,echo,head,tail,cat) and apply
some attention to a small subset of the {\color{verbcolor}{\verb#bash#}}
command shell's capabilities. 

We will manipulate file names using the bash shell's variable
substitution rules together with looping and conditional constructs.

We will introduce several ways to automatically edit files using:
{\color{verbcolor}{\verb#sed,awk and vi#}}. 

In short, we want to master a small subset of the powerful Unix environment
to make brief one-liners (Bon Mots!) to get the job done quickly.

\subsection{Unix and its One-Line Editors}

Unix commands expect its input from a special file known as STDIN
(the command line or a file that is piped or otherwise redirected
INTO a command). It does its work and sends its output to STDOUT
which may be the terminal's display or to a pipe or file. If there
are any errors they are sent to the special file known as STDERR.
STDIN,STDOUT and STDERR are three main default streams through which
data (information) flows. 

\clearpage
\subsubsection{Sed (the stream editor)}

The {\color{verbcolor}{\verb#sed#}} program accepts its
input from STDIN does work specified by one of a few command
line expressions and sends its output to STDOUT.

In tIRAF task {\color{verbcolor}{\verb#crmedian#}} is an odd-ball
in that it refuses to work on more than one file at a time.

\textbf{\emph{Strategy:}} To make crmedian work on all files, we will use Unix
to create a {\color{verbcolor}{\verb#cl#}} file on the fly, and
run that file. This violates so many principles! But it works.

{\color{verbcolor}
\begingroup \fontsize{10pt}{10pt}
\selectfont
%%\begin{Verbatim} [commandchars=\\\{\}]
\begin{verbatim} 
! ls -1 *fts *fits *fit > l.l
! echo "crmedian.unlearn" > crall.cl
! echo "crmedian.sigma    = ''" >>crall.cl
! echo "crmedian.residual = ''" >>crall.cl
! echo "crmedian.crmask   = ''" >>crall.cl
! echo "crmedian.median   = ''" >>crall.cl
! cat l.l | sed  -e 's/\(^.*$\)/crmedian \1 c_\1/' >>crall.cl
cl < crall.cl
\end{verbatim}
\endgroup
%% \end{Verbatim}
}

\textbf{\emph{Prose:}} \textbf{\emph{Hey}}, PyRAF, send a few lines to
bash. \textbf{\emph{Hey}}, bash, first make a list of all the fits
files. Then use {\color{verbcolor}{\verb#echo#}} to send some text (a
cl command to set an internal variable related to crmedian)
{\color{verbcolor}{\verb#echo "crmedian.unlearn"#}} via redirection
using a single greater-than-sign to start a new script file
{\color{verbcolor}{\verb#> crall.cl#}}. Then send a few similar lines,
using double greater-than-signs to append text to the emerging new
script file. (These lines tell iraf to ignore making mask files!.)
Then, bash, run the {\color{verbcolor}{\verb#cat#}} command to read
the list of filenames {\color{verbcolor}{\verb#l.l#}} and send (pipe)
{\color{verbcolor}{\verb#|#}} the filenames (one at a time) to the
stream editor {\color{verbcolor}{\verb#sed#}}.

\textbf{\emph{Hey}}, {\color{verbcolor}{\verb#sed#}}, use a single
command ({\color{verbcolor}{\verb#-e#}}) of
{\color{verbcolor}{\verb#'s/\(^.*$\)/crmedian \1 c_\1/'#}} that uses a
{\color{verbcolor}{\verb#g/re/p#}} like command. Globally
({\color{verbcolor}{\verb#g#}}) use a regular expression
({\color{verbcolor}{\verb#re#}}) {\color{verbcolor}{\verb#\(^.*$\)#}}
to do substitute the contents of the line with
{\color{verbcolor}{\verb#crmedian \1 c_\1#}} -- and followed by a
print ({\color{verbcolor}{\verb#p#}}).  The regular expression is of
the form {\color{verbcolor}{\verb#\(#}} start a group.  Into that
group accept all characters satisfying
{\color{verbcolor}{\verb#^.*$#}} (from the start of the line
{\color{verbcolor}{\verb#^#}} all characters
{\color{verbcolor}{\verb#.*#}} stopping at the end of the line
{\color{verbcolor}{\verb#$#}}. Then close the group
{\color{verbcolor}{\verb#\)#}}. Send the output to STDOUT using the
double greater-than-signs, of the form of the command
{\color{verbcolor}{\verb#crmedian#}} the original filename held in
group 1 {\color{verbcolor}{\verb#\1#}} and create the output filename
by prepending a {\color{verbcolor}{\verb#c_#}} to the original
filename held in group 1 {\color{verbcolor}{\verb#c_\1#}}.

Now, hehe --- \textbf{\emph{Hey}}, cl (yes you are cl script that is already running so
start a new one), and take that file that {\color{verbcolor}{\verb#sed#}}
just made as the commands to run!


\clearpage
\subsection{A bit about Unix output redirection}

\index{Unix!io redirection}
Bash redirection symbols:

{\color{verbcolor}
\begin{quote}
\begingroup \fontsize{10pt}{10pt}
\selectfont
%%\begin{Verbatim} [commandchars=\\\{\}]
\begin{verbatim}
command >  file              -- Create file, redirect STDOUT into that file.
command >> file              -- Append output to file.
command |  command  >file    -- Connect command1 output to command2 input.
command >  file              -- Regular output into file, errors to console.
command >2 file              -- Regular output to console, errors to file.
command >  file2>&1          -- Both output and errors to file.
\end{verbatim}
\endgroup
%% \end{Verbatim}
\end{quote}
}


Example of Unix using the ``plumbing'' paradigm to ``redirect'' data ``flows'' from
one command to another.

\begin{quote}
{\color{verbcolor}{\verb#!ls -1 F*lat*fits | tee in.txt | sed -e 's/^/out_/' > out.txt#}} \\
{\color{verbcolor}{\verb#cat in.txt#}} \\
{\color{verbcolor}{\verb#cat out.txt#}}
\end{quote}

Not a single interactive editor is used.

\textbf{\emph{Tasks:}} {\color{verbcolor}{\verb#ls, tee, sed#}} and
 {\color{verbcolor}{\verb#cat#}}. The bang ({\color{verbcolor}{\verb#!#}})
has PyRAF send the rest of the line (all the Unix part) to the bash
shell.

\textbf{\emph{Prose:}} {\color{verbcolor}{\verb#ls#}} ``lists'' files
according to the wild card supplied and the output stream ``flows''
via the Unix redirection operator the vertical bar
({\color{verbcolor}{\verb#|#}}). The flow ``ties'' (plumbs?) the
output from {\color{verbcolor}{\verb#ls#}} to the input of
{\color{verbcolor}{\verb#tee#}}; where tee makes a copy of the stream
as it flows along sending one copy into a file called in.txt (mirror
of the ls output) and then the other along the stream to another pipe
redirector that passes the stream into {\color{verbcolor}{\verb#sed#}}; where a ``regular
expression'' is applied to each line. The regular expression says for
the start of each line prepend a {\color{verbcolor}{\verb#out_#}}
sub-string with sed's output being redirected into a file called
{\color{verbcolor}{\verb#out.txt#}}. The {\color{verbcolor}{\verb#cat#}} command
provides a peek at the output for our approval, because, wait for it, I always
check my work.

\clearpage
\subsection{The ``vi'' visual editor in batch mode}

The Unix editor, {\color{verbcolor}{\verb#vi#}} can be run by
supplying a series of commands {\color{verbcolor}{\verb#-c#}} and
ending with commands to {\color{verbcolor}{\verb#write#}} and
{\color{verbcolor}{\verb#quit#}} the process.

E.g.: Same crmedian example:

{\color{verbcolor}
\begingroup \fontsize{10pt}{10pt}
\selectfont
%%\begin{Verbatim} [commandchars=\\\{\}]
\begin{verbatim} 
!ls -1 *fits > l.l
!vi -es -c "%s/\(^.*$\)/crmedian \1 c_\1/" -c "w" -c "q" l.l
# here l.l is the lines crmedian filename.fits c_filename.fits
#but the preamble is not there.
! echo "crmedian.unlearn" > crall.cl
! echo "crmedian.sigma    = ''" >>crall.cl
! echo "crmedian.residual = ''" >>crall.cl
! echo "crmedian.crmask   = ''" >>crall.cl
! echo "crmedian.median   = ''" >>crall.cl
! cat l.l >> crall.cl
cl < crall.cl
\end{verbatim}
\endgroup
%% \end{Verbatim}
}

\textbf{\emph{Prose:}} Hey, {\color{verbcolor}{\verb#ls#}}, make that list of
filenames.  Now, {\color{verbcolor}{\verb#vi#}} change the
{\color{verbcolor}{\verb#l.l#}} file -- no copies! So use the echo
sequence, as above, and start the {\color{verbcolor}{\verb#crall.cl#}}
file. Then copy the {\color{verbcolor}{\verb#ci#}} modified contents
to the {\color{verbcolor}{\verb#crall.cl#}} file. Then, cl -- do your
trick.
\clearpage
\subsection{AWK -- a very powerful editor/language}

AWK is a very powerful programming environment in-and-of itself.
Here is the crmedian problem

{\color{verbcolor}
\begingroup \fontsize{10pt}{10pt}
\selectfont
%%\begin{Verbatim} [commandchars=\\\{\}]
\begin{verbatim} 
!ls -1 *fits > l.l
# here l.l is the lines crmedian filename.fits c_filename.fits
#but the preamble is not there.
! echo "crmedian.unlearn" > crall.cl
! echo "crmedian.sigma    = ''" >>crall.cl
! echo "crmedian.residual = ''" >>crall.cl
! echo "crmedian.crmask   = ''" >>crall.cl
! echo "crmedian.median   = ''" >>crall.cl
! cat l.l | awk  '/./ {print "crmedian $0 c_$0";}' >> crall.cl
cl < crall.cl
\end{verbatim}
\endgroup
%% \end{Verbatim}
} 

\textbf{\emph{Prose:}} Hey PyRAF, do the usual
{\color{verbcolor}{\verb#ls -1 *fits > l.l#}} trick and then use a
one-liner with {\color{verbcolor}{\verb#awk#}} to write the commands
into place. So, {\color{verbcolor}{\verb#awk#}} take each line that
matches the pattern {\color{verbcolor}{\verb#/./#}} (at least
something on the line) and print the whole line
{\color{verbcolor}{\verb#$0#}} as the input filename and a
{\color{verbcolor}{\verb#c_$0#}} as the output filename.


The entire script generation could be managed differently
with an ``awk'' script:

{\color{verbcolor}
\begingroup \fontsize{10pt}{10pt}
\selectfont
%%\begin{Verbatim} [commandchars=\\\{\}]
\begin{verbatim} 
#!/bin/awk
# say this is awk_crmedian
BEGIN{
   print "crmedian.unlearn";
   print "crmedian.sigma    = ''";
   print "crmedian.residual = ''";
   print "crmedian.crmask   = ''";
   print "crmedian.median   = ''";
}
/./ {print "crmedian $0 c_$0";}
\end{verbatim}
\endgroup
%% \end{Verbatim}
}

\textbf{\emph{Prose:}} Hey, we wrote our own script {\color{verbcolor}{\verb#awk_crmedian#}}
and made it executable and put in into {\color{verbcolor}{\verb#~/iraf#}}.


Then...

{\color{verbcolor}
\begingroup \fontsize{10pt}{10pt}
\selectfont
%%\begin{Verbatim} [commandchars=\\\{\}]
\begin{verbatim} 
!ls -1 *fits | awk_crmedian > crall.cl
cl < crall.cl
\end{verbatim}
\endgroup
%% \end{Verbatim}
}

\textbf{\emph{Prose:}} Hey, PyRAF have bash do the usual
{\color{verbcolor}{\verb#ls -1 *fits#}} trick, and pipe the output
into our script called {\color{verbcolor}{\verb#awk_crmedian#}}; hey,
{\color{verbcolor}{\verb#awk_crmedian#}} dump the output to our hacked
script crall.cl; then PyRAF use the cl interpreter to run that script.

\clearpage
\section{Headers should be fixed}

NASA listing of popular keywords.
\url{https://fits.gsfc.nasa.gov/fits_dictionary.html}


The critical header to fix is IMAGETYP. IRAF wants the
values {\color{verbcolor}{\verb#zero,dark,flat,object#}} all
in \textbf{\emph{lower case}}!. Some observatories will
write a {\color{verbcolor}{\verb#NaN#}} (not a number) in
WCS values that cause some tasks to literally blow up. \index{IMAGETYP!values}

\textbf{\emph{Note:}} these files are usually in a {\color{verbcolor}{\verb#ccddb#}}
directory where the file types could be bent to the local usage.


Different observatory engineering teams and various camera vendor's
software packages (\cite{IRAFMotherLoad,SBFITSEXT}) add headers
that are not 'consistent' with what IRAF wants.

Other IRAF keywords are: DATE-OBS, EXPTIME, GAIN, RDNOISE and
IMAGETYP. \index{IRAF!very important keywords} The specification
really wants a {\color{verbcolor}{\verb#Z#}} at the end of the
DATE-OBS to be perfectly clear the string is a ``Zulu'' time
(UTC). UTC is the best to use as good records allow corrections.

WCS keywords may be inaccurate, and different observatories and
software packages write these headers. 



 

See Table \ref{table:MainKeywords} for the main keywords we're about to discuss.

We will \textbf{\emph{add}}, \textbf{\emph{translate}} and \textbf{\emph{remove}} keywords.

\vspace{-.15cm}
\begin{enumerate}[resume]\addtolength{\itemsep}{-0.5\baselineskip}
   \item   Header KEYWORDS can not be easily changed -- E-GAIN can not be directly changed to GAIN.
\begin{enumerate}\addtolength{\itemsep}{-0.5\baselineskip}
   \item   Add the right one
   \item   Delete the older bad one
\end{enumerate}
   \item   Value can be changed.
   \item   Missing keywords are simply added.
\end{enumerate}



\textbf{\emph{BTW:}} Add a filter of ``none'' to zeros and darks:

{\color{verbcolor}{\verb#hselect *fits $I "(IMAGETYP ?= 'dark' || IMAGETYP ?= 'zero')" > l.l#}} \\
{\color{verbcolor}{\verb#hedit @l.l FILTER none add+ update+ del- ver- show-#}}

\textbf{\emph{And}}, strip WCS keywords from zeros and darks, or at
least guarantee no ``NaN'' or ``inf'' values are in WCS headers for
non-object files. This requires listing a few files representative of
the WCS imposed for the night's observations; and creating a series of
hedits to meet needs:

{\color{verbcolor}{\verb#hedit wildcard Comma-Keyword-List del+ ver- show- update+#}}

\begin{table}[h!]
%\phantomsection
%\addcontentsline{toc}{section}{ TOC CAPTION}
% \setlength{\belowcaptionskip}{6pt} % adjust space under caption abovecaptionskip
% \renewcommand{\arraystretch}{1.3} % adjust line spacing
%\small{
%\begin{minipage}{\textwidth}     % for footnotes in table.
%\caption[TOC]{Main Keywords}
\centering
\begin{tabular}{ l  l l}
%\MakeShortVerb{\|}
%\multicolumn{n}{fmt}{text for merged cols}
\hline
Right      & Wrong                                          & The Fix  \\
\hline
GAIN       & EGAIN                                          & remove the E \\
RDNOISE    & usually missing                                & use obsutil findgain task \\
IMAGETYP   & IMGTYPE  EXPTYPE                               & hedit   \\
DATE-OBS   & DATE and TIME                                  & merge, add 'Z' timezone part for UTC \\
OBJECT     & Mostly right                                   & remove spaces?   \\
FILTER     & ties to {\color{verbcolor}{\verb#ccdred$#}}... & to taste   \\
\hline
PIXSIZE1   & in degrees                                     & translate \\
PIXSIZE2   & in degrees                                     & translate  \\
CCDSUM     & <int> or ( x   y)  pixels                      & translate \\
\hline
OBSGEO-B   & LATITUDE SITE-LAT                              & add/translate \\
OBSGEO-L   & LONGITUD SITE-LONG                             & add/translate \\
OBSGEO-H   & ALTITUDE  (Greisen  +)                         & add/translate \\
LONGPOLE   & (usually missing)                              & add $\equiv$ 0 Earth \\
LATPOLE    & (usually missing)                              & add $\equiv$ 0 Earth\\
\hline
OBSERVER   & Main observer (team name)                      & add \\
OBSERV01   & names of the observers...                      & add \\
OBSERV02   & names of the observers...                      & add \\
\hline
OBJRA      & Sexagesimal                                    & add/translate \\
OBJDEC     & Sexagesimal                                    & add/translate \\
\hline
SATURATE   & float                                          & sextractor \\
\hline
%%\DeleteShortVerb{|}
\end{tabular}
%%\end{minipage}    %% for footnotes  r@{.}l
\caption{Main Keywords for IRAF}
\label{table:MainKeywords}
%%} % end small etc
\end{table}


IRAF wants \textbf{\emph{lower-case}} text as the IMAGETYP keyword's
value.\footnote{It matters. See files like
  \dhl{iraf/noao/imred/ccdred/ccddb/kpno/camera.dat}.}

\begin{table}[h!]
\centering
\begin{tabular}{ l  l  l }
\hline
IMAGETYP         & Flexberry Synonym   & Meaning            \\
\hline
\multicolumn{3}{c}{KPNO Vocabulary}                         \\
\hline
OBJECT (0)           & object          &                    \\
DARK (1)             & dark            &                    \\
PROJECTOR FLAT (2)   & flat            &                    \\
SKY FLAT (3)         & other           &                    \\
COMPARISON LAMP (4)  & other           &                    \\
BIAS (5)             & zero            &                    \\
DOME FLAT (6)        & flat            &                    \\
\hline
\multicolumn{3}{c}{Maxim/DL Vocabulary}                     \\
\hline
\multicolumn{3}{c}{Advanced FlexBerry Vocabulary}           \\
object               & object          & Target exposure    \\
light                & object          &                    \\
target               & object          &                    \\
sci                  & object          &                    \\
science              & object          &                    \\
light frame          & object          &                    \\
bias                 & zero            &                    \\
bias frame           & zero            &                    \\
zero                 & zero            & Zero/Bias          \\
dark                 & dark            & Dark               \\
dark frame           & dark            & Dark               \\
flat field           & flat            & Flat (Package dependent) \\
flat                 & flat            &                    \\
comp                 & comp            & Spectro Comparison 1D \\
comp2d               & comp2d          & Spectro Comparison Image 2D \\
1d                   & 1d              & 1D spectrum naxis=1 \\
(unknown)            & unknown         & unknown/missing keyword \\
%% ones-based: \cline{a-b}
\hline
\end{tabular}
\caption{Image Type Keywords}
\label{table:ImageTypeKeywords}
\end{table}
\clearpage
\subsection{Fix Things Up}

Introducing the IRAF commands:  

{\color{verbcolor}{\verb#imheader#}}    \\
{\color{verbcolor}{\verb#hselect#}} and \\
{\color{verbcolor}{\verb#hedit#}}

\subsubsection{imheader}
\index{commands!imheader}
Look at an ``image'' ``header'':

\begin{quote}
{\color{verbcolor}{\verb#imheader filelist#}} \\
{\color{verbcolor}{\verb#imheader filelist long+ user+ | less#}}
\end{quote}

{\color{verbcolor}{\verb#imheader filelist#}} reports a very basic
list. It blows up on MEF\footnote{FITS Multi-Extension-FITS.} files!
Easy way to see a combine operation made a MEF.

\begin{quote}
{\color{verbcolor}{\verb#imheader filelist long+ user+#}}
\end{quote}

will show all the user headers.

E.g.:

\begin{quote}
{\color{verbcolor}{\verb#imhead c_IC1396_WC_0010.fits#}}
\end{quote}

reports:

{\color{darkgreen}
\begin{quote}
\begingroup \fontsize{10pt}{10pt}
\selectfont
%%\begin{Verbatim} [commandchars=\\\{\}]
\begin{verbatim}
imhead c_IC1396_WC_0010.fits
c_IC1396_WC_0010.fits[1374,1099][ushort]: IC1396_WC
\end{verbatim}
\endgroup
%% \end{Verbatim}
\end{quote}
}
The template for the output:

\begin{quote}
{\color{darkgreen}{\verb#filename[NAXIS1,NAXIS2][16-bit integer]: OBJECT#}}
\end{quote}

E.g.:

\begin{quote}
{\color{verbcolor}{\verb#imhead c_IC1396_WC_0010.fits long+#}}
\end{quote}

reports:

{\color{darkgreen}
\begin{quote}
\begingroup \fontsize{8pt}{8pt}
\selectfont
%%\begin{Verbatim} [commandchars=\\\{\}]
\begin{verbatim}
c_IC1396_WC_0010.fits[1374,1099][ushort]: IC1396_WC
No bad pixels, min=0., max=0. (old)
Line storage mode, physdim [1374,1099], length of user area 2673 s.u.
Created Tue 16:28:49 02-Oct-2018, Last modified Tue 16:28:57 02-Oct-2018
Pixel file "c_IC1396_WC_0010.fits" [ok]
\end{verbatim}
\endgroup
%% \end{Verbatim}
\end{quote}
}
E.g.:

\begin{quote}
{\color{verbcolor}{\verb#imhead c_IC1396_WC_0010.fits long+ user+#}}
\end{quote}

reports:
{\color{darkgreen}
\begin{quote}
\begingroup \fontsize{8pt}{8pt}
\selectfont
%%\begin{Verbatim} [commandchars=\\\{\}]
\begin{verbatim}
imhead c_IC1396_WC_0010.fits long+ user+
c_IC1396_WC_0010.fits[1374,1099][ushort]: IC1396_WC
No bad pixels, min=0., max=0. (old)
Line storage mode, physdim [1374,1099], length of user area 2673 s.u.
Created Tue 16:28:49 02-Oct-2018, Last modified Tue 16:28:57 02-Oct-2018
Pixel file "c_IC1396_WC_0010.fits" [ok]
EXTEND  =                    F / File may contain extensions
BSCALE  =           1.000000E0 / REAL = TAPE*BSCALE + BZERO
BZERO   =           3.276800E4 /
ORIGIN  = 'NOAO-IRAF FITS Image Kernel July 2003' / FITS file originator
DATE    = '2018-10-02T21:07:29' / Date FITS file was generated
IRAF-TLM= '2018-10-02T22:28:57' / Time of last modification
OBJECT  = 'IC1396_WC'          / Name of the object observed
DATE-OBS= '2018-06-27T10:54:26' /YYYY-MM-DDThh:mm:ss observation start, UT
\end{verbatim}
\endgroup
%% \end{Verbatim}
\end{quote}
}
\clearpage
\subsubsection{IRAF COMMAND hselect}
\index{commands!hselect} 

The {\color{verbcolor}{\verb#hselect#}} command takes filenames from
a list of files, answers with a list of keyword values that satisfy a
boolean expression.

\begin{quote}
{\color{verbcolor}{\verb#hselect filelist Comma-Keyword-List boolean#}}
\end{quote}

The special keyword {\color{verbcolor}{\verb#$I#}} stands in for the
related image's filename. The other keywords as they appear in the header.

Its the boolean expression that is the main power. A simple
{\color{verbcolor}{\verb#yes#}} means to accept all files.

The boolean {\color{verbcolor}{\verb#"(IMAGETYP ?= 'Bias')"#}} looks
at all files, only acts on those where the {\color{verbcolor}{\verb#IMAGETYP#}}
keyword's value roughly matches or {\color{verbcolor}{\verb#looks like (?=)#}}
the string. \textbf{\emph{Note:}} The expressions starts with double-quotes to allow
use of single-quotes on the inside of the boolean test. Note: use {\color{verbcolor}{\verb#==#}}
to precisely match the entire quoted string, and {\color{verbcolor}{\verb#!=#}} to
``not'' ``match'' the entire quoted string. This follows the ``C'' bash string comparison
paradigm.

Other handy boolean tests:

\begin{quote}
\begingroup \fontsize{10pt}{10pt}
\selectfont
%%\begin{Verbatim} [commandchars=\\\{\}]
\begin{verbatim}
Line: Expression
1  hselect c_Flat_0011.fits $I,NAXIS1,NAXIS2 yes
2  hselect *fits   $I "(IMAGETYP ?= 'Bias')"   > l.l
3  hselect *fits   $I "(IMAGETYP ?= 'Dark')"   > l.l
4  hselect *fits   $I "(IMAGETYP ?= 'Flat')"   > l.l
5  hselect *fits   $I "(IMAGETYP ?= 'Light')"  > l.l
6  hselect *fits   $I,OBSGEO-B,OBSGEO-L,OBSGEO-H  yes
7  hselect *fits   $I,LONGPOLE,LATPOLE         yes
8  hselect *fits   $I,FILTER,EXPTIME,IMAGETYP "(IMAGETYP == 'object' || IMAGETYP == 'flat')"
9  hselect c_*fits $I,EXPTIME,FILTER "(IMAGETYP == 'flat')"
10 hselect c_*fits $I "(IMAGETYP == 'flat' && FILTER == 'HAlpha' && EXPTIME==25 )" > l.l
11 hselect c_*fits $I "(IMAGETYP == 'dark' && EXPTIME=1800 )" > l.l
\end{verbatim}
\endgroup
%% \end{Verbatim}
\end{quote}

In the above (prose):

\vspace{-.15cm}
\begin{enumerate}[resume]\addtolength{\itemsep}{-0.5\baselineskip}
   \item   For the file {\color{verbcolor}{\verb#c_Flat_0011.fits#}} show the name
and size (NAXIS1 and NAXIX2)
   \item   Find all files where IMAGETYPE ``looks like'' Bias, and list only the
filename {\color{verbcolor}{\verb#($I)#}} one-file-per-line into a
temp file called {\color{verbcolor}{\verb#l.l#}}.
(next step! {\color{verbcolor}{\verb#hedit @l.l IMAGETYP zero add+ ver- show- update+#}}
(next next step!) {\color{verbcolor}\\
{\verb#imcombine @l.l zmaster.fits combine=mode#}}
In other words, don't depend on the file's name to be right and while we're at it
change the header value to the IRAF default, then might as well cook up the
master zero file.
   \item   Make a list of all the dark files, next step? (hedit...)
   \item   Make a list of all the flat files
   \item   Make a list of  all the possible science objects
   \item   Look at the site location for all the files (more to very all were updated)
   \item   Look at the LONGPOLE and LATPOLE values, make sure they are right
   \item   OK, object and flat files need to me ``zero subtracted'', so get the list
the follow with a \\
{\color{verbcolor}{\verb#imarith @l.l - zmaster.fits z_@l.l#}}
   \item   List all the names of all the flats.
   \item   Make a list: test for IMAGETYP of flat, FILTER of HAlpha, and an EXPTIME of 25s
and save to temp file l.l. This type of list can be used with a \\
{\color{verbcolor}{\verb#imcombine @l.l flatmaster.fits ...#}} .
   \item   For the SBO ABG cameras, where we have a lot of 1800s Darks, make a list of them
\end{enumerate}
\clearpage
\subsubsection{IRAF COMMAND hedit}
\index{commands!hedit}
\begin{quote}
{\color{verbcolor}{\verb#hedit listoffiles keyword newvalue switches#}}
\end{quote}

Examples of hedits, for all files with a wildcard {\color{verbcolor}{\verb#*fits#}}
or from a list we made with {\color{verbcolor}{\verb#files#}} or with
{\color{verbcolor}{\verb#hselect#}}:

\begin{quote}
\begingroup \fontsize{8pt}{8pt}
\selectfont
%%\begin{Verbatim} [commandchars=\\\{\}]
\begin{verbatim}
1  hedit *fits OBSERVER 'teamwoody'               add+ ver- show- update+
2  hedit @l.l IMAGETYP zero                       add+ ver- show- update+
3  hedit *fits PIXSIZE1 "(@'XPIXSZ')"             add+ ver- show- update+
4  hedit *fits PIXSIZE2 "(@'YPIXSZ')"             add+ ver- show- update+
5  hedit *fits RDNOISE 5.31                       add+ ver- show- update+
6  hedit *fits GAIN 0.26                          add+ ver- show- update+

# how to delete a lot of headers in one go (bad wcs for example)
7  hedit @l.l CTYPE1,CRVAL1,CRPIX1,CDELT1,CROTA1  del+ update+ show- ver-
\end{verbatim}
\endgroup
%% \end{Verbatim}
\end{quote}

\vspace{-.15cm}
\begin{enumerate}[resume]\addtolength{\itemsep}{-0.5\baselineskip}
   \item   Add/change the OBSERVER keyword to be our team name.
   \item   For the list of {\color{verbcolor}{\verb#IMAGETYP ?= 'Bias'#}} we
got from a {\color{verbcolor}{\verb#hselect#}}, change to the proper
{\color{verbcolor}{\verb#zero#}} string value.
   \item Change a keyword. Can't do that, but we can add a new keyword
     with the old one's value. The
     {\color{verbcolor}{\verb#"(@'XPIXSZ')"#}} picks up the value of
     an existing {\color{verbcolor}{\verb#XPIXSZ#}} keyword and uses
     it as the value for the new keyword. Hint:\\
     {\color{verbcolor}{\verb#hedit @l.l XPIXSIZ del+ update+ show- ver-#}}.
   \item   Same for the matching PIXSIZE2 \menu YPIXSZ
   \item Run findgain, or use a calculator, and determine the real
     GAIN and RDNOISE. Here add/update RDNOISE.
   \item   Add/update GAIN
   \item WCS solutions are often not good. There are several lines of
     the associated keywords. Here is one hedit of several to strip
     the WCS from the file using the
     {\color{verbcolor}{\verb#delete#}} switch.
\end{enumerate}


Example of using hselect and hedit together:

{\color{verbcolor}
\begin{quote}
\begingroup \fontsize{10pt}{10pt}
\selectfont
%%\begin{Verbatim} [commandchars=\\\{\}]
\begin{verbatim}
hselect *fits $I "(IMAGETYP ?= 'Bias')" > l.l
cat l.l
hedit @l.l IMAGETYP zero   add+ ver- show- update+
\end{verbatim}
\endgroup
%% \end{Verbatim}
\end{quote}
}
The line {\color{verbcolor}{\verb#cat l.l#}} shows us the list,
and lets us check it is what we thought we asked for with the
{\color{verbcolor}{\verb#hselect#}}.

\subsubsection{Binning}

Binning is all messed up with vendor software. The IRAF keywords
dictionary cites: {\color{verbcolor}{\verb#CCSSUM = 2#}} or
{\color{verbcolor}{\verb#CCDSUM = (3 1)#}} (no comma) for
spectroscopy.

\begin{quote}
{\color{verbcolor}{\verb#hedit *fits CCDSUM  "((@'XBINNINB)"#}}
\end{quote}

works for us with symmetric binning and standard imaging.

Handling \gls{noise} is important.


% from QuickCrib.


For observations: create a desktop link (shortcut) of
\dhl{\textasciitilde{}/Desktop/Today} folder, and initialize all the
acquisition software to save files into that directory. At the end of
the night simply rename/move the directory to your prefered location
and immediately create a new \dhl{\textasciitilde{}/Desktop/Today}
folder. \dhl{Hint}: means several packages are confugured once, and may
be quickly deployed when taking data.

The observation date should be the \dhl{date of sunset} at the observatory. This
avoids midnight rolling over to the next date and really confusing where
pieces of the night winds up.

Directory Names: I prefer ddMMMyyyy where MMM is the three letter
abbreviated name of the month. This absolutely avoids ambiguity
between cultures. Using a ISO-like format of YYYYmmdd or YYYY-mm-dd,
is OK but it is to hard to read amongst the files, and diminishes
\dhl{tab completion}.

A night of observing benefits from a decent plan. The plan focuses on the target
(science) stars, should proceed from west to east, and once the meridian is reached,
the targets should be observed as they come onto the meridian. This is not always
possible, try to catch targets as close to the meridian as possible. This reduces
the airmass for the targets and where possible the reference stars.

A plan should match target (science) stars with reference stars. Reference stars should be
as close as possible to the target star as possible to maintain the same air mass
and to mitigate flexure in the system. 

\subsection{File Names}

In Unix, all files are based on the machine's single ``root''
directory (/). Directories on other machines can be "remote" mounted
using any of SMB or NFS features of Unix. For Windows mounts, use
SMB. Note: SMB is more susceptable to malware.

Never use spaces, punctuation, plus/minus. \\
Do use underscore or your own shorthand: plus (p) and minus (m) is not a bad idea.

Start with the target star's program name, or try to use HD Catalog
numbers where possible.  If not HD, then fall back to HIP, after that,
you are on your own. I prefer to use SIMBAD MAIN\_ID names, as they
are easier to lookup while observing.

Follow the Target observation with an \_Cal (calibration )or \_Ref
(reference star) ... etc. This will help sub-sort the file names together
for a target.

In any event get a simple convention, and stick with it. Try to provide
the name of the stars in the filename.

It is bad form to put necessary information in the filename that is
not in the header for a file. So update the fits files with these
additional information as soon as possible. The IRAF \dhl{nhedit} (new
hedit) command inside of a cl file is a great way to do this --
discussed herein.

Keep a hand-written paper log!

\subsubsection{Spectroscopy Example}


A spectrographic observation involves a matched pair of target and reference star
exposures. The operations are:
\vspace{-.15cm}
\begin{enumerate}\addtolength{\itemsep}{-0.5\baselineskip}
%\setcounter{enumi}{N}
   \item   Slew to the science target
\vspace{-.15cm}
\begin{enumerate}\addtolength{\itemsep}{-0.5\baselineskip}
%\setcounter{enumi}{N}
   \item   Take three exposures, prevent loss of data when cosmic rays hit.
   \item   BEFORE SLEW, take three cal lamp images. 
\end{enumerate}

   \item   Slew to the reference star
\vspace{-.15cm}
\begin{enumerate}\addtolength{\itemsep}{-0.5\baselineskip}
%\setcounter{enumi}{N}
   \item   Take three exposures, prevent loss of data when cosmic rays hit.
   \item   BEFORE SLEW, take three cal lamp images. 
\end{enumerate}
\end{enumerate}

The reference star should be chosen to be as close to the target as possible. In
modern observing settings the composition of the atmosphere may vary. Care should
be taken to avoid reference stars with extinction -- unless the extinction matches
to a reasonable degree to the target.

If the position of the target and reference fall onto different parts of the slit,
a spectrographic flat process needs to be imployed. Again, three flat exposures
with each pointing.

Most acquisition packages do not adaquately recored all information for the observation.
These may be added as soon as possible after the raw data is acquired. It is important
to have the DATE-OBS, RA, and DEC for all pointings to record the orientation of
the instrument package. Yes, this includes calibration lamp and flat images. It does
not apply to zero or dark exposures.

The SIMBAD name \index{SIMBAD!name} contains variable numbers of spaces and other
decorations like plus/minus signs, dashes, comas, asterisks, and a handful of other
characters. This makes easy string comparisons difficult. If possible use the same
text as the SIMBAD, perhaps in a separate KEYWORD in the fits file.

Using a database helps to manage the particulars of observing. The
PostgreSQL database offers the best suite of tools for the task, is
free, well documented, and maintained. Unlike MySQL or MariaDB --
PostgreSQL, a database is very isolated and does not share data
between other databases. PostgreSQL uses a database subdivision called
a ``schema'' that allows data to be partitioned as needed. If you are
used to MySQL treat two schema within one database in the rough/same
context as two databases.

The package Q3C (Section \ref{sec:PostgreSQL}) has additional details of
special tools for matching stellar coordinates. PostgreSQL supports programming
special functions in Python. 

\clearpage

\begin{table}[h!]
\centering
\begin{tabular}{| l | l | l |}
\hline
KEYWORD  & psql field & type  \\
\hline
RA       & ra         &  float,   \\ 
DEC      & dec        &  float,   \\ 
OBJECT   & targetname &  text,    \\ 
SIMBAD   & simbadname &  text,    \\ 
VMAG     & vmag       &  float,   \\ 
SPTYPE   & sptype     &  text,    \\ 
EQUINOX  & equinox    &  float -- EPOCH depricated,   \\ 
RADESYS  & radesy     &  text, -- ICRS,FK4, J2000  \\ 
DATE-OBS & dateobs    &  text,  \\ 
         & obtypes    &  json, -- keyval   \\ 
         & fluxes     &  json,  -- keyval   \\ 
\hline
\multicolumn{3}{| c |}{ Rotator}  \\
\hline
ROTANG   & rotang     &   \\
\hline
\multicolumn{3}{| c |}{ Instrument}  \\
\hline
SLITANG  & slitang  \\
SLIT     & slit       &  float,  \\
SLITWIDE & slitwid    &  float, -- [arcseconds]  \\
GRATING  & grating    &  float, -- [l/mm]  \\
\hline
\multicolumn{3}{| c |}{ Telescope}  \\
\hline
FOCALEN  & focalen    & float, -- focal length in mm  \\
APERTURE & aperture   & float, -- clear aperture  \\
\hline
\end{tabular}
\caption[FITS HEADER]{Necessary FITS Keywords. Autoguider from ING}
\label{table:NecessaryFITSKeywords}
\end{table}


\begin{table}[h!]
\centering
\begin{tabular}{| l | l | l |}
\hline
\multicolumn{3}{| c |}{Autoguider} \\
\hline
KEYWORD  & psql field & type  \\
\hline
AUTACQT   & AUTACQT  & float,      --  Integration time for acquisition (s)      \\
AUTGUIT   & AUTGUIT  & float,      --  Integration time for guiding (s)          \\
AUTWINDO  & AUTWINDO & integer     --  Autoguider window used                    \\
AUTMODE   & AUTMODE  & text,       --  Autoguider mode used (e.g. Cass, UES)     \\
AUTXSEE   & AUTXSEE  & float,      --  Mean FWHM in X for guide star (arcsec)    \\
AUTYSEE   & AUTYSEE  & float,      --  Mean FWHM in Y for guide star (arcsec)    \\
AUTXSRMS  & AUTXSRMS & float,      --  RMS on FWHM in X for guide star (arcsec)  \\
AUTYSRMS  & AUTYSRMS & float,      --  RMS on FWHM in Y for guide star (arcsec)  \\
AUTMAGNI  & AUTMAGNI & float,      --  Mean instrumental magnitude of guide star \\
AUTTRANS  & AUTTRANS & float,      --  RMS variation of magnitude of guide star  \\
AUTMAGOF  & AUTMAGOF & float,      --  Magnitude zero-point                      \\
AUTSTATN  & AUTSTATN & integer     --  No of data-points used for stats. calculations  \\
\hline
\multicolumn{3}{| c |}{Camera}                                                   \\
\hline
BIN1     & bin1       & binning naxis1                                           \\
BIN2     & bin2       &  binning naxis2                                          \\
CCDSUM   & ccdsum     & text default '(1 1)', -- binning 1 and 2                 \\
PIXSCAL1 & PIXSCAL1   & float default 3.74, -- arcsec on axis  i                 \\
PIXSCAL2 & PIXSCAL2   & float default 3.74, -- arcsec on axis  i                 \\
DETECTOR & detector   & string default 'Sony IMX571', -- chip name               \\
\hline
\multicolumn{3}{| c |}{Weather}                                                  \\
\hline
AMBIENT             & ambient   & float  , -- ambient temperature                \\
PRESSURE            & pressure  & float  , -- ambient pressure mmHG              \\
RELHUMID            & relhumid  & float  , -- percent                            \\
%%  ones-based: \cline{a-b}
\hline
\end{tabular}
\caption[FITS HEADER]{Necessary FITS Keywords. Autoguider from ING}
\label{table:NecessaryFITSKeywords}
\end{table}

\ltodo{FLEXSPEC KW}{Add keywords for flexspec.}
\begingroup \fontsize{10pt}{10pt}
\selectfont
\begin{verbatim} 
# from ING
CENWAVE = REAL               / Approx central wavelength (A)      *
LINESMM = REAL               / Number of lines/mm of the grating  *
DISPERSI= REAL               / Nominal dispersion (A/mm)          *
ISIANAMO= CHAR               / Position of anamorphotic lens
ISIHWAVE= CHAR               / Position of half-wave polarizer
ISIQWAVE= CHAR               / Position of quarter-wave polarizer
ISIDEKKE= CHAR               / Position of dekker
ISISLITU= CHAR               / Slit unit (long, multi or fibre)
ISISLITW= INTEGER            / Slit width (microns)
PAOFFSET= REAL               / Slit p.a. - tel. p.a. (degs)
ISIFCP  = CHAR               / Field-lens, calcite, polaroid or clear
ISIBFOLD= CHAR               / Position of blue fold
ISIRFOLD= CHAR               / Position of red fold
ISISARM = CHAR               / Data from red arm, blue arm or FOS ?
ISIFILTA= CHAR               / Name of filter A in ISIS arm      *
ISIFILTB= CHAR               / Name of filter B in ISIS arm      *
ISIHART = CHAR               / Configuration of Hartmann shutter *
ISICOLL = INTEGER            / Position of collimator (microns)  *
ISIGRAT = CHAR               / Name of grating                   *
ISITHETA= INTEGER            / Angle of grating (mdegs)          *
ISIBXDIS= CHAR               / Position of blue x disperser
ISITIGHT= CHAR               / Light tight ?
ISISDOOR= CHAR               / Status of slit door
ISICELL = CHAR               / State of grating-cell clamps      *
ISIGDOOR= CHAR               / Status of grating door
ISILID  = CHAR               / Status of ISIS lid
PAOFFSET= REAL               / Slit p.a. - telescope p.a. (degrees)
\end{verbatim}
\endgroup


The keywords RA, DEC, EQUINOX ,RADESYS, LATITUDE, LONGITUD, ALTITUDE,
together with DATE-OBS, EXPTIME and JD plays into the calculation of
airmass. There is no real physics equation that solves a true value
for airmass as the atmospheric conditions are handled in different
ways between authors. The critical thing is feed algorithms with
proper values for the existing IRAF algorithms in use.

The target list has the name of the desired target, vmag, exptime and the
name of the associated reference star (also name, vmag, exptime). The
reference star is associated with a 

All files are presumed to have been trimmed, zero, dark and flat corrected.
For CMOS images the dark step should be omitted. PyRAF will be used
together with the pyraflogin3.py helper file. In the case of the sci2d.fits
file, one would expect to see \verb=f_d_z_t_scied.fits= as the file name
to match in order trimming, zero subtraction, dark subtraction and flat
correction. 

To simplify this section's typesetting, these filenames will be used:

\begin{table}[h!]
%\phantomsection
%\addcontentsline{toc}{section}{ TOC CAPTION}
% \setlength{\belowcaptionskip}{6pt} % adjust space under caption abovecaptionskip
% \renewcommand{\arraystretch}{1.3} % adjust line spacing
%\small{
%\begin{minipage}{.8\textwidth}     % for footnotes in table.
%\caption[TOC]{Filenames and meanings for the example.}
\centering
\begin{tabular}{| l | l |}
%\MakeShortVerb{\|}
%\multicolumn{n}{fmt}{text for merged cols}
\hline
Filename          &  Description                                   \\ 
\hline
sci2d.fits        &  the science target                            \\ 
comp2d.fits       &  comparison lamps at science target pointing   \\ 
flat2d.fits       &  flat lamps at science target pointing    \\ 
sci\_ref2d.fits   &  target's A0V reference star                   \\ 
comp\_ref2d.fits  &  target's A0V reference star                   \\ 
flat\_ref2d.fits  &  target's A0V reference star                   \\ 
sci.fits          &  1D files for the target star                  \\ 
comp.fits         &  1D files for the target's cal                 \\ 
flat.fits         &  1D files for the tarrget's flat               \\ 
sci\_ref.fits     &  1D files for the reference's star             \\ 
comp\_ref.fits    &  1D files for the reference's cal              \\ 
flat\_ref.fits    &  1D files for the reference's flat             \\ 
%% ones-based: \cline{a-b}
\hline
%%\DeleteShortVerb{|}
\end{tabular}
%%\end{minipage}    %% for footnotes  r@{.}l 
\caption[TOC]{Filenames and meanings for the example.}
\label{table:Filenamesandmeaningsfortheexample.}
%%} % end small etc
\end{table}

The Processing steps are:

\begin{table}[h!]
%\phantomsection
%\addcontentsline{toc}{section}{ TOC CAPTION}
% \setlength{\belowcaptionskip}{6pt} % adjust space under caption abovecaptionskip
% \renewcommand{\arraystretch}{1.3} % adjust line spacing
%\small{
%\begin{minipage}{.8\textwidth}     % for footnotes in table.
%\caption[TOC]{Turse Summary of the steps.}
\centering

Steps to be taken.

\begin{tabular}{| l | l |}
%\MakeShortVerb{\|}
%\multicolumn{n}{fmt}{text for merged cols}
\hline
Task     & Summary    \\ 
\hline
apall    & apextract the trace from sci2d.fits to make sci.fits.    \\ 
apall    & Using sci apertures, extract comp2d.fits to make the comp.fits.    \\ 
Identify & comp.fits to create the wavelength data.    \\ 
dispcor  & comp aperture with the sci.fits image, produce a wavelength calibrated fits file.    \\ 

%% ones-based: \cline{a-b}
\hline
%%\DeleteShortVerb{|}
\end{tabular}
%%\end{minipage}    %% for footnotes  r@{.}l 
\caption[TOC]{Turse Summary of the steps.}
\label{table:TurseSummaryofthesteps.}
%%} % end small etc
\end{table}

\subsection{Observatory}

The file \dhl{/opt/iraf/lib/iraf/noao/lib/obsdb.dat} needs an entry:
\begingroup \fontsize{10pt}{10pt}
\selectfont
\begin{verbatim} 
observatory "elsauce"
    name = "Obstech Observatory"
    longitude = -70.76498
    latitude  = -30.47085
    altitude  = 1564
    timezone  = 4
\end{verbatim}
\endgroup

Note: Some pinhead put unicode into the file, it may need to be repaired.
